\section{Introduction}
\label{sec:intro}

\subsection{Longevity risk is priced off the interval}

A life annuity is a bet on a distribution. An annuity book's best estimate
is a functional of the central mortality projection; its risk margin, its
longevity capital and the price at which a longevity swap or bulk annuity
transfer clears are functionals of its \emph{spread}: under the
Solvency~II standard formula a permanent 20\% fall in rates
\citep{eiopa2015delegated}, under an internal model the 99.5\% one-year
quantile of the change in own funds. The balance sheet carries a tail
of a forecast distribution, whose adequacy rests on one property: a nominal
95\% interval must contain the outcome about 95\% of the time. Too narrow
under-reserves, too wide over-prices the hedge, and neither shows in a
point-accuracy comparison. Classical mortality models and their neural
extensions all emit intervals
\citep{lee1992modeling,marino2023neural,miyata2022extending,schnurch2022point};
whether those are \emph{calibrated} is a separate question from whether they
exist.

\subsection{The break that makes the question answerable}

Until 2019 the high-income Human Mortality Database (HMD) mortality
surfaces improved smoothly enough that almost any trend
extrapolation looked calibrated in sample. The years 2020--2024 broke that
surface: life expectancy at birth fell from 2019 to 2020 in 27 of the 29
countries analysed by \citet{aburto2022quantifying}. Intervals built from
data ending in 2019 thus met their first genuine out-of-sample test of the
modern record, unscored
until now: the closest study, \citet{barigou2023bayesian}, evaluates by
leave-future-out validation, log score and CRPS but simulates its pandemic,
real post-shock data not existing in 2021. It exists now: the HMD vintage of
15~June~2026 is the first with final single-age data through 2024 for the
twenty populations audited here.

\subsection{What the literature measures}

Across a full-text search of about twenty modelling papers
(Section~\ref{sec:related}), point accuracy is the default and often the
only metric: \citet{richman2021neural} and
\citep{scognamiglio2022calibrating,nigri2019deep,petnehazi2019mortality,perla2021time}
report mean squared error and no interval, coverage or proper score;
elsewhere intervals are drawn but never checked against realisations
\citep{miyata2022extending}. Where coverage is checked it carries no
sharpness penalty, under pass rules in which over-coverage is fulfilment
\citep{schnurch2022point} or a coverage probability of $1$ is best
\citep{marino2023neural}; the latter also reports without follow-up that
nominal 95\% Lee--Carter ARIMA intervals cover ``around 50\%'' of realisations
and, for Spanish
males, ``stable around 33\%'', over a window containing no break, and
\citet{miao2025gradient} call the closest of marginal coverages 90\%, 94\%
and 96\% ``the most properly calibrated'', again with no proper score and no
width penalty. \citet{dowd2010backtesting} do backtest densities ex post,
but on one population and without a break. Absent, so far as we can find:
PIT or rank histograms, measured joint (path) coverage, formal
forecast-comparison testing but for \citet{shang2018model}, and---in the
2025--2026 conformal mortality line
\citep{shang2025intervals,shang2026conformal,shang2026subnational}---any
crossing of a designated break. The governing principle is old: maximise
sharpness subject to calibration
\citep{gneiting2007probabilistic,gneiting2007strictly}. Coverage without
width is not calibration, nor width without coverage sharpness; we score
both.

\subsection{This study}

Do stochastic, machine-learning and neural mortality models attain nominal
interval coverage across the COVID-19 break; does that depend on the
family, the mechanism or both; and does it differ when nothing breaks? We
treat \emph{model family} and \emph{uncertainty mechanism} as separate crossed
factors---ten families, seven mechanisms,
fifty admissible cells---in three regimes: a \emph{stable} control
(expanding origins 1990--2014, every test year before 2020), the
\emph{shift} regime (train to 2019, test 2020--2024), and a \emph{placebo}
break (train to 1913, test 1914--1922) testing whether the finding is
specific to COVID-19. Population counts differ by regime and family---twenty
across the break and sixteen for the Gaussian process, eighteen and sixteen
in the stable control, eleven in the placebo---so are stated wherever a mean
is taken.

Five hypotheses were fixed and hash-stamped before any model touched real
data, and are reproduced verbatim from the pre-registration.

\begin{hypothesis}[H1]
Model rankings by RMSE and by CRPS disagree (rank correlation $<1$; at least
one top-3 inversion) in the shift regime.
\end{hypothesis}
\begin{hypothesis}[H2]
Nominal 95\% intervals under-cover in the shift regime for every model
family; the shortfall exceeds the stable-regime shortfall.
\end{hypothesis}
\begin{hypothesis}[H3]
Joint path coverage over $h = 1,\dots,5$ is materially below marginal
coverage at the same nominal level for every family (reported gap; test:
paired difference $>0$).
\end{hypothesis}
\begin{hypothesis}[H4]
Coverage failure is non-uniform in age: shift-regime coverage at ages 65--99
is lower than at ages 25--64 for every family.
\end{hypothesis}
\begin{hypothesis}[H5]
Interval miscalibration propagates to derived quantities: empirical coverage
of 95\% intervals for $e_{65}$ and annuity factor $\annuity$ (2\%) departs
from nominal by more than the corresponding stable-regime departure.
\end{hypothesis}

No neural-versus-classical ordering is registered; the crossed design
estimates it.

\subsection{Contributions}

This paper proposes no new mortality model. Its contributions are:
\begin{enumerate}
\item \emph{A reusable audit protocol.} The crossed design is implemented by
  a harness pre-registered and hash-stamped before the first fit, so a
  coverage failure can be attributed to the architecture, the mechanism or
  their interaction. Each of the fifty cells emits pathwise predictive
  samples through one \texttt{fit}/\texttt{predict\_dist} interface and is
  scored through one code path, so no model can be flattered by a bespoke
  metric and joint path coverage is computable. Two gates passed
  before any real-data result was read: nominal coverage for a correctly
  specified model on synthetic truth from a known Poisson Lee--Carter
  process, and numerical parity with R \pkg{StMoMo} on the same HMD subset.
  Claims are contrasts within admissible sub-grids, never main effects over
  a ragged grid; inference is a wild cluster bootstrap over populations, a
  pandemic being one shock and not twenty. It transfers to future breaks.
\item \emph{The audit, and the finding the control forces.} The first
  evaluation of probabilistic mortality forecasts---classical, neural and
  conformal---against realised 2020--2024 mortality, every interval built
  strictly from pre-2020 data, beside the stable control. The Gaussian
  process's forty-year training floor removes four populations, and registered
  Addendum~4 caps its window at sixty years
  (Table~S11). Most of the miscalibration predates the pandemic: the native
  95\% intervals of the classical cores already cover 0.702 (Lee--Carter),
  0.682 (Poisson Lee--Carter) and 0.731 (Renshaw--Haberman) of
  outcomes in the control, against 0.692, 0.624 and 0.654 across the break
  (Table~S3). The break widens a pre-existing failure rather than creating
  one; the shift regime alone would have misattributed all of it to
  COVID-19. CBD is where the break contributes most, and only the control
  shows it: 0.901 in the quiet decades against 0.799 across 2020--2024,
  about two-thirds of its shift-regime departure break-induced and a third
  not. \Hone\ holds, the registered top-three inversion present, with RMSE
  and CRPS orderings correlating at $\rho = 0.95$ and RMSE with the Poisson
  log score at $\rho = 0.12$ (Table~S2).
\item \emph{Where the failure sits.} Not on a neural-versus-classical split:
  across the break the negative-binomial head (0.933) and the
  CNN under dropout (0.962) hold marginal coverage near nominal and the
  Gaussian process reaches 0.905, while the LSTM (0.647 ensembled, no better
  in the control at 0.638) and the ensembled neural Lee--Carter (0.694) sit
  with the classical cores far below nominal in both regimes (Table~S3). Nor
  is the mechanism contrast a pandemic effect: on the 95\% interval score
  $\IS$ the split conformal wrapper beats the native interval for the same
  four classical families in both regimes, same sign, different sizes
  (Poisson Lee--Carter $+1.038$ to $+1.895$), significantly for all four
  across the break ($p$ from 0.004 to 0.018) but in the control only for
  Lee--Carter and Poisson Lee--Carter ($p = 0.006$ and $p = 0.032$, against
  $p = 0.086$ for Renshaw--Haberman and $p = 0.198$ for CBD). The native
  interval beats its wrapper, significantly in both regimes, for the three
  families whose native law was right: the negative-binomial head
  ($-0.873$ control, $-1.024$ shift), the Gaussian process ($-0.772$,
  $-0.413$) and the sparse VAR ($-0.236$, $-0.318$), all at $p \le 0.0002$
  in the control (Diebold--Mariano with a wild cluster bootstrap over
  populations; Table~S9). Conformal wrapping repairs families whose native
  law was wrong and damages those whose was right, in quiet periods as much
  as in the pandemic. The 1914--1922 placebo does not reproduce that
  ordering: the contrast reverses for the four classical families, none
  significantly ($p$ from 0.071 to 0.689), while the
  native advantage persists for the negative-binomial head ($-0.461$,
  $p = 0.008$) and the Gaussian process ($-0.172$, indistinguishable from
  zero).
\item \emph{Registered directions that fail, reported as such.} \Htwo\
  registered one-sided under-coverage across the break for \emph{every}
  family, worsening on the control. The registered two-sided departure
  $\lvert c - 0.95 \rvert$ grows from stable to shift in 27 of the 50
  primary cells and raw coverage falls in 34 of 50 (Table~S3), but a fall is
  degradation only for an arm that began below nominal, and most arms above
  nominal move \emph{towards} it (sparse VAR native 0.988 to 0.937). The
  pandemic mostly sorts families by the side of nominal they started on---a
  tendency, not a law, describing 33 of the 50 cells. Five
  over-covering arms widen their departure instead: four copula-conformal
  bands, of which the Poisson Lee--Carter's crosses nominal on the way down
  (0.960 to 0.939, its registered departure growing from 0.010 to 0.011, the
  one exception), and the CNN under dropout, above nominal in both regimes
  (0.954 to 0.962). Twelve under-covering arms move back towards nominal,
  eight conformal wrappers recalibrated on pre-2020 residuals (EnbPI on
  Lee--Carter, 0.890 to 0.915), the rest LSTM arms already too low to
  fall further. \Hfour\ fails in its universal form and settles a
  disagreement with the prior: where \citet{dowd2010backtesting} found
  age~84 better behaved than age~65, the native Lee--Carter interval covers
  0.888, 0.760 and 0.501 across ages 0--24, 25--64 and 65--99 in the stable
  control, before any break (Table~S5). Old-age concentration is real for
  the Lee--Carter line and the LSTM arms but not universal: the CBD arms
  cover worse below age 65 than above it (native, shift: 0.626 against
  0.849). CBD is fitted on ages 55--99, so that 0.626 is a 55--64 stub
  rather than the registered 25--64 band and is not comparable with the
  full-age families.
\item \emph{Resolution of the failure by horizon and in money.} Joint path
  coverage over $h = 1,\dots,5$ lies below marginal coverage in all 100 of
  the 100 (regime, cell) pairs, with mean gaps of 0.127 stable and 0.168
  across the break (Table~S4). That count is largely arithmetic: independent
  horizons alone would put joint coverage at $0.95^{5} = 0.774$, a gap of
  0.176 wider than either mean, so the registered comparison
  cannot on its own evidence miscalibration. Against the
  informative benchmark---each arm's joint coverage under independence
  \emph{at its own marginal rate}---joint coverage exceeds that value in 50
  of 50 stable cells, 50 of 50 shift cells and 45 of 46 placebo cells, by as
  much as 0.34: horizons are strongly positively dependent, a dependence
  \citet{dowd2010backtesting} name but never measure and \Hthree\ here
  quantifies. The registered re-specified comparator, each cell's
  model-implied joint coverage (Addendum~3, \S8), is not computable from the
  columns the production runner emitted, so \Hthree\ carries no registered
  verdict here; what survives is that joint coverage falls below the nominal
  independence value 0.774 in 21 of the 50 stable cells and 30 of 50 across
  the break. The copula conformal band, the only mechanism calibrated for
  the path, narrows the gap furthest on average without
  closing it (0.805--0.927 joint coverage across the break), but not
  everywhere: on the negative-binomial head the dropout arm's gap of 0.075
  beats its 0.085. Derived quantities are then far worse calibrated than the
  rates beneath them in every arm but three with
  computable derived coverage---the two CBD arms, integrated on ages 55--99
  alone, and the sparse
  VAR's native law (Section~\ref{sec:actuarial})---and this too is a
  quiet-period property: in the stable control the native Lee--Carter 95\%
  interval covers 0.253 of realised annuity factors $\annuity$ and 0.256 of
  realised $e_{65}$, and the LSTM ensemble 0.157 of annuity factors.
  \Hfive's registered comparative clause---a larger departure from nominal
  across the break than in the control---holds for 14 of the 20
  distributional arms on $e_{65}$ and 12 of 20 on $\annuity$, the exceptions
  being arms already so far below nominal in the control that the break
  could only move them towards it (Table~S6).
\end{enumerate}
